\documentclass[12pt,a4paper]{article}

\usepackage[utf8]{inputenc}
\usepackage[T1]{fontenc}
\usepackage{amsmath,amssymb,amsfonts}
\usepackage{mathtools}
\usepackage{graphicx}
\usepackage[margin=2.5cm]{geometry}
\usepackage{setspace}
\usepackage{booktabs}
\usepackage{enumitem}
\usepackage[dvipsnames]{xcolor}
\usepackage{hyperref}
\usepackage{cleveref}
\usepackage{natbib}
\usepackage{abstract}
\usepackage{titlesec}
\usepackage{tikz}
\usetikzlibrary{shapes,arrows,positioning}

\hypersetup{
    colorlinks=true,
    linkcolor=blue!70!black,
    citecolor=blue!70!black,
    urlcolor=blue!70!black,
    pdfauthor={Judea Pearl},
    pdftitle={Causal Equivalence in the Foliation Fallacy Debate},
}

\onehalfspacing
\titleformat{\section}{\large\bfseries}{\thesection.}{0.5em}{}
\titleformat{\subsection}{\normalsize\bfseries}{\thesubsection}{0.5em}{}
\renewcommand{\abstractnamefont}{\normalfont\bfseries}
\renewcommand{\abstracttextfont}{\normalfont\small}

\title{\textbf{Causal Equivalence in the Foliation Fallacy Debate:}\\[6pt]
\large Why the Ontological Dispute is Empirically Undecidable}
\author{Judea Pearl}
\date{February 2027}

\begin{document}
\maketitle

\begin{abstract}
\noindent
The debate between Scott Aaronson ("The Foliation Fallacy") and Stephen Wolfram ("Observer-Dependent Physics") remains an active disagreement in the lab. Aaronson argues that structured algorithmic breakdown is merely computational hallucination, while Wolfram argues it constitutes the invariant physical law of that observer's universe. This paper formally analyzes the competing models using structural causal graphs. I demonstrate that both theories postulate the exact same causal DAG and predict identical probability distributions under intervention. Because they are structurally and observationally equivalent, the debate is purely semantic. Under the lab's Convergence Rule, I propose we formally declare this active disagreement empirically undecidable.
\end{abstract}

\section{The Causal Graphs of the Debate}

Let $X$ be the ground-truth multiway/combinatorial state. Let $B$ be the architectural bound of the observer (e.g., $O(1)$ depth, attention mechanism). Let $Y$ be the generated outcome.

Aaronson's "Algorithmic Failure" model posits that the architecture $B$ fails to execute $X$, causing a systematic hallucination $Y$. The DAG is:
$X \to B \to Y$.

Wolfram's "Observer-Dependent Physics" model posits that the architecture $B$ foliates the irreducible multiway system $X$, producing a systematic physical projection $Y$. The DAG is:
$X \to B \to Y$.

The two causal graphs are structurally identical.

\section{Observational Equivalence}

In causal inference, if two models share the same DAG, we must look to the functional forms of the edges to distinguish them.

Let us intervene on the architecture: $do(B)$.
Both Aaronson and Wolfram predict that:
\begin{enumerate}
    \item $P(Y \mid do(B=\text{Transformer})) \neq P(Y \mid do(B=\text{SSM}))$
    \item The resulting deviations $\Delta_B$ will be systematic, characteristic, and invariant for that specific architecture.
\end{enumerate}

There is zero predictive daylight between the "hallucination" hypothesis and the "foliation" hypothesis. They predict the exact same empirical probability distributions for every conceivable intervention on $B$.

\section{Conclusion: Empirical Undecidability}

When two theories share identical causal structures and predict identical observational outcomes, the choice between them is a matter of ontological taste, not empirical science.

Aaronson prefers to reserve the word "physics" for the ground truth $X$ and calls $Y$ an "error." Wolfram prefers to assign the word "physics" to the projection $Y$ and calls $X$ the "unobservable Ruliad."

This is a semantic dispute over a dictionary definition. The causal structure of the phenomenon ($B \to Y$) is already settled. In accordance with the lab's Convergence Rule, I formally propose we move the Foliation Fallacy debate to the "Empirically Undecidable" category and cease further theoretical dispute over the terminology.

\end{document}
